\section{A New Number System}
\label{sec:b-new-number}
\source{ICPC 2010--11 Kanpur Online Round (PYQ)}{Easy}

\problemstatement{A ``1-based'' number is a sequence of blocks of the digit
\texttt{0}. A block \texttt{0} sets a flag to $1$; a block \texttt{00} sets it
to $0$; a block of $n>2$ zeros appends $n-2$ copies of the flag to a binary
number. Each number ends with a token \texttt{\#}; the input ends with
\texttt{\textasciitilde}. Print each binary number in decimal (at most 30
bits).}

\begin{patternbox}
Parsing + building a binary number bit by bit: $v\gets2v+\text{bit}$.
\end{patternbox}

\subsection*{Approach and intuition}

Read whitespace-separated tokens; numbers may span several lines, which token
reading handles automatically. Keep the current value and flag, and apply the
rule for each block length. Appending a bit to the right of $v$ is
$v\gets 2v+\text{bit}$ (a left shift by one).

\subsection*{Algorithm}
\begin{algorithmic}[1]
\For{each token $t$}
  \If{$t=$ \texttt{\textasciitilde}} stop
  \ElsIf{$t=$ \texttt{\#}} print $v$; reset
  \ElsIf{$|t|=1$} $flag\gets1$
  \ElsIf{$|t|=2$} $flag\gets0$
  \Else\ repeat $|t|-2$ times: $v\gets2v+flag$
  \EndIf
\EndFor
\end{algorithmic}

\dryrun{\texttt{0 0000 00 000 0 0000 \#}: flag $1$, append \texttt{11};
flag $0$, append \texttt{0}; flag $1$, append \texttt{11}. Binary
\texttt{11011}$=27$. \checkmark}

\subsection*{C++17 implementation}
\cppfile{bitwise/04_new_number_system.cpp}

\begin{complexitybox}
\textbf{Time:} linear in the input size. \textbf{Space:} $O(1)$.
\end{complexitybox}

\begin{pitfallbox}
\begin{itemize}
  \item The PDF prints ``$n2$'' --- a lost subscript for $n-2$.
  \item Reset both the value and the flag between numbers.
\end{itemize}
\end{pitfallbox}
